% CHAPTER 6 - CONCLUSION AND FUTURE WORK
\chapter{Conclusion and Future Work}

\section{Conclusion}
This project designed, implemented and evaluated the Context Memory Framework, a
context-aware Natural Language to SQL assistant built around a persistent, governed
context-memory layer. The work started from four practical shortcomings of plain NL2SQL in
an organization: business knowledge missing from the schema, stateless interaction, schemas
too large for a prompt, and the risk of executing model-written SQL.

CMF answers each of them with a concrete mechanism. Business knowledge is stored as
semantic, procedural and episodic memories with scope, visibility, importance, validity and
provenance; model-extracted memories stay pending until a person approves them, every edit
is versioned and audited, corrections supersede old definitions, and one-time overrides are
never stored. Statelessness is addressed by working memory: recent turns and a structured
record of the previous query, with follow-up questions retrieved using the previous
question. Large schemas are reduced to the most relevant tables plus their foreign-key
neighbours. Safety is enforced deterministically: one read in a rolled-back transaction with
a row cap and masking, or one single-table change that an admin must approve and that is
re-validated before it runs.

Retrieval combines dense and sparse search with reciprocal rank fusion, a weighted
relevance score and a cross-encoder reranker, always after authorization filters, and every
decision is recorded and shown in the Context Inspector. The evaluation showed that all 66
automated tests pass and that, on a 120-case synthetic benchmark, the current pipeline
reaches Recall@3 of 0.79 and MRR of 0.73, includes the target context in 97.2\% of the
checked cases, keeps definitions current in all superseding cases and leaks no memory
across users or organizations. These results show that the context layer reliably delivers
the right, authorized knowledge to the model; they do not yet measure how much this
improves the correctness of the generated SQL.

\section{Limitations}
The following limitations of the current prototype are known and documented:
\begin{itemize}[leftmargin=*, topsep=4pt, itemsep=3pt]
    \item \textbf{No real authentication.} Identity is the organization and user identifier
          sent with each request, \texttt{default-org}/\texttt{default-user} are created
          automatically as owner, and the role chosen on the landing page is a
          demonstration, not a security boundary.
    \item \textbf{Credential storage.} Connection passwords are base64-encoded, not
          encrypted.
    \item \textbf{Single LLM provider.} Only Google Gemini is supported, through the
          \texttt{google-generativeai} package, which is now deprecated. Several settings
          (\texttt{LLM\_PROVIDER}, \texttt{ENABLE\_LOCAL\_LLM}, \texttt{MAX\_RETRIES},
          \texttt{ENFORCE\_READ\_ONLY\_SQL}, \texttt{MEMORY\_DEFAULT\_STATUS}) exist in the
          configuration but are not read by the code.
    \item \textbf{Latency.} A question can cause up to three LLM calls (SQL, answer and
          background extraction), and the embedding and reranking models load on first use.
    \item \textbf{Scalability.} Without Qdrant, retrieval scans every authorized memory row.
          Embeddings are stored twice (JSON in the platform database and in Qdrant). Token
          counts are approximated as characters divided by four.
    \item \textbf{Operations.} CORS allows all origins, schema changes use a small list of
          \texttt{ALTER} statements instead of a migration tool, and the MySQL connection URL
          uses the \texttt{pymysql} driver, which is not listed in \texttt{requirements.txt}.
    \item \textbf{Evaluation scope.} The benchmark is synthetic and does not execute an LLM;
          end-to-end SQL accuracy was not measured.
\end{itemize}

\section{Future Scope}

\subsection{End-to-End Evaluation}
The most important next step is to measure SQL execution accuracy with and without the
memory layer, for example on BIRD \cite{li2023bird}, whose external-knowledge evidence could
be loaded as memories instead of being given per question, and on the five sample
databases with hand-written questions and expected results.

\subsection{Security Hardening}
Real authentication (for example OpenID Connect) should replace the identifier-based
identity, connection secrets should be encrypted, CORS should be restricted, and
row-level policies could be added to the access model alongside table-level grants.

\subsection{Retrieval and Memory Improvements}
The scoring weights and the reranker weight were set by hand; they could be tuned on a
labelled set of real questions. Memory consolidation could merge related memories, and
schema drift could be handled by using the stored \texttt{schema\_hash} to flag memories that
refer to dropped or renamed columns. An exact tokenizer would make the token budget precise.

\subsection{Provider Flexibility, Performance and Usability}
The LLM client could be migrated to the maintained Google SDK and extended to other
providers, including local models, through the existing \texttt{LLM\_PROVIDER} setting;
embedding and reranking models could be preloaded at start-up to remove first-request
latency. Rendering Markdown in answers and charting numeric results would make the
interface clearer for non-technical users.
